\section{第\ref{sec:chemoutput}章　化学出力}

心臓への2本の配線とその速さの違い、ホルモンカスケードの負帰還、分量の頻度による符号、概日発振器とその光による同期は直接測定されている。境界$K<8$は本章で導いた制御理論の定理である。カスケードのフィードバックそのものが脈動を生むというのは仮説であり、それを強く支持する実験がある。

{\footnotesize
\setlength{\tabcolsep}{3pt}\renewcommand{\arraystretch}{1.15}
\begin{longtable}{>{\raggedright}p{0.5cm}>{\raggedright}p{4.0cm}>{\raggedright}p{5.6cm}>{\raggedright}p{2.3cm}>{\raggedright\arraybackslash}p{1.7cm}}
\toprule
№ & 主張 & 実験と測定内容 & 出典 & 状態 \\
\midrule\endfirsthead
\toprule
№ & 主張 & 実験と測定内容 & 出典 & 状態 \\
\midrule\endhead
1 & 心臓のペースメーカーは単独で毎分約$100$回拍動する。安静時はブレーキが踏まれており、心拍数はふつう$60$--$70$程度 & ヒト、心臓への2本の配線をともに完全に遮断：固有心拍数は安静時心拍数より高く、加齢とともに下がり、成人で毎分約$100$拍。したがって安静時はブレーキが優位である。安静時心拍数$60$--$70$は典型値で、個別には引用していない。 & \cite{JoseCollison1970ihr} & 測定済み \\
2 & アクセル\nt{NORA}とブレーキ\nt{ACET}はローパスフィルタで、ブレーキの方が速い~\eqref{eq:gasbrake} & イヌ、心臓の迷走神経または交感神経を周波数変調刺激し、心房レートへの伝達関数を測定：どちらの経路もローパスフィルタで、交感神経側は遮断周波数がはるかに低く、約$1.7\,\text{s}$の純粋な遅れが加わる。特性は平均緊張度に依存するので、線形式~\eqref{eq:gasbrake}は近似である。 & \cite{Berger1989} & 測定済み \\
3 & ブレーキが速いのは、\nt{ACET}がセカンドメッセンジャーなしでカリウムチャネルを開き、\nt{NORA}は反応の連鎖を通じて作用するから & 心房細胞の膜パッチ：アセチルコリンで開くカリウムチャネルは、受容体と共役したGタンパク質の$\beta\gamma$サブユニットが、細胞内のセカンドメッセンジャーを介さずに開く。cAMPを介する\nt{NORA}の経路はここでは引用していない。 & \cite{Logothetis1987gbg} & 測定済み \\
4 & 血液はおよそ1分で全身を一周する。血中の\nt{ADRE}は受容体をもつすべての臓器に届く & 一般的な桁の見積もり。本章でもそのように述べており、出典は引用していない。 & --- & 推定 \\
5 & ホルモンカスケードは負帰還で囲まれている：最終段のホルモンが前段を抑える & 実験の総説：副腎皮質ホルモンは、下垂体からのACTH放出と視床下部からの放出ホルモン放出を、速い、中間、遅いの三つの時間域で抑制する。 & \cite{KellerWood1984fb} & 測定済み \\
6 & 同一の3段は$K<8$で安定~\eqref{eq:cascadestable}、境界での周期は$2\pi\tau/\sqrt3\approx3.6\tau$ & 本章で導出：周波数$\sqrt3/\tau$で3段は位相を$180^\circ$ずらし、$8$分の1に減衰させる。 & 本章で導出 & 理論 \\
7 & レベルの誤差は$1/(1+K)$で、この調節器はおよそ$10\,\%$より高い精度を保てない（命題~\ref{prop:cascade}） & 比例フィードバックに対する6行目の帰結。実際の軸は複数の段と複数の時間域にフィードバックをもつ（5行目）ので、この境界はモデル~\eqref{eq:cascade}についてのものであり、測定値ではない。 & 本章で導出 & 理論 \\
8 & $K=4$と$7$ではレベルが落ち着き、$K=12$では周期$3.7$--$3.9\,\tau$の規則正しい脈動（図~\ref{fig:cascade}） & \texttt{models/\allowbreak chem\_\allowbreak models.py}による計算：$K=12$で連続する周期は$3.9$、$3.7$、$3.8$、$3.7\,\tau$。$K=4$では計算終了時の振幅$0.003$。 & \texttt{models/\allowbreak chem\_\allowbreak models.py} & モデル \\
9 & ストレスホルモンは約1時間に1回のパルスで放出される。パルスはフィードバックそのものが生み、独立した発振器はない & ラット、放出ホルモンの持続注入：ACTHとコルチコステロンは生理的なほぼ1時間の頻度で振動する。パルスに視床下部からのリズム入力は必要ない。遅れのあるフィードバックをもつ下垂体--副腎モデルがこの振動を再現する。 & \cite{Walker2012glc, Walker2010} & 仮説 \\
10 & 受け手はレベルではなく分量の頻度を読む & 性腺刺激ホルモン放出ホルモンを自ら分泌できないサル：持続注入では下垂体ホルモンの放出は回復せず、1時間に1回の分量投与で回復する。持続注入に戻すと応答は再び消える。ハイパスフィルタとしての受容体の疲労は解釈である。 & \cite{Belchetz1978} & 測定済み \\
11 & 分量の頻度が違えば、同じ腺の異なる細胞への作用も違う & 同じサル：分量を1時間に$2$--$5$回に増やすと下垂体ホルモンが両方とも下がる。$3$時間に1回に減らすと一方（LH）が下がり他方（FSH）が上がるので、両者の比は頻度で決まる。 & \cite{Wildt1981gnrh} & 測定済み \\
12 & 概日リズムは数時間の遅れをもつ細胞内の負帰還が生む & 総説：時計遺伝子のタンパク質が遅れて自らの遺伝子の転写を抑える。転写と翻訳のループが約1日の周期を生む。 & \cite{TakahashiJS2017} & 測定済み \\
13 & 主発振器は数万個のニューロンからなる群で、体の他の部分を同期させる & 総説：視交叉上核が主概日発振器である。培養した個々のニューロンは独立した発振器で、ネットワークがそれらを同期させる。マウスでは片側に約$10^4$個のニューロン。 & \cite{Welsh2010scn} & 測定済み \\
14 & ヒトの固有周期は$24.18$時間 & 薄暗い光のもと、1日から切り離したスケジュールで暮らすヒト：メラトニン、体温、コルチゾールのリズムの周期は若年者と高齢者で平均$24.18$時間、ばらつきは小さい。 & \cite{Czeisler1999} & 測定済み \\
15 & 光は位相同期ループのように発振器を合わせる~\eqref{eq:pll}。必要な補正は1日約$11$~min & ヒト、1日のさまざまな時刻に$6.7$時間の明るい光を1回照射：振幅$5$時間のタイプ1位相反応曲線。体温最低点より前では遅れ、後では進む。式~\eqref{eq:pll}は標準モデル、$11$~min $=0.18$時間は算術である。 & \cite{Khalsa2003prc} & 測定済み \\
16 & 光がなければ同期はなく、リズムは固有周期で進む & 光をまったく知覚できない全盲者で、通常の社会生活を送る人：約半数でメラトニンとコルチゾールのリズムが1日と一致しない固有周期で進む。 & \cite{Sack1992blind} & 測定済み \\
\bottomrule
\end{longtable}}
